\documentclass{article}
\usepackage{amsmath,amssymb,amsthm}
\usepackage{algorithm}
\usepackage[noend]{algpseudocode}
\usepackage{enumitem}
\usepackage{hyperref}
\usepackage{bm}
\newtheorem{theorem}{Theorem}
\newtheorem{lemma}{Lemma}
\newtheorem{assumption}{Assumption}
\newtheorem{definition}{Definition}
\newtheorem{corollary}{Corollary}

\begin{document}

\section*{Algorithm Name}
\textbf{Proximal Gradient Method for Non‑Convex Composite Optimization (PGM‑NC)}  

We consider the composite problem  
\[
\min_{x\in\mathbb{R}^d}\;\Phi(x)\;\stackrel{\text{def}}{=}\;f(x)+g(x),
\tag{1}
\]
where $f$ is smooth (possibly non‑convex) and $g$ is proper, lower‑semicontinuous (possibly non‑smooth).  

The iteration reads
\[
x_{k+1}\;=\;\operatorname{prox}_{\eta g}\!\bigl(x_k-\eta\nabla f(x_k)\bigr),
\qquad k=0,1,\dots,
\tag{2}
\]
with stepsize $\eta>0$.

---------------------------------------------------------------------

\section*{Mathematical Formulation}
\begin{itemize}
    \item \textbf{Proximal operator:}
    \[
    \operatorname{prox}_{\eta g}(v)\;\stackrel{\text{def}}{=}\;
    \arg\min_{u\in\mathbb{R}^d}\Bigl\{ g(u)+\frac{1}{2\eta}\|u-v\|^{2}\Bigr\}.
    \tag{3}
    \]
    \item \textbf{Gradient mapping (stationarity measure):}
    \[
    G_{\eta}(x_k)\;\stackrel{\text{def}}{=}\;\frac{1}{\eta}\bigl(x_k-x_{k+1}\bigr)
    =\frac{1}{\eta}\Bigl(x_k-\operatorname{prox}_{\eta g}\!\bigl(x_k-\eta\nabla f(x_k)\bigr)\Bigr).
    \tag{4}
    \]
\end{itemize}

---------------------------------------------------------------------

\section*{Pseudocode}
\begin{algorithm}[H]
\caption{Proximal Gradient Method for Non‑Convex Composite Optimization}
\begin{algorithmic}[1]
\State \textbf{Input:} initial point $x_0\in\mathbb{R}^d$, stepsize $\eta>0$, max iterations $K$.
\For{$k=0$ \textbf{to} $K-1$}
    \State Compute gradient $g_k \gets \nabla f(x_k)$.
    \State Take proximal step 
    \[
    x_{k+1}\gets \operatorname{prox}_{\eta g}\bigl(x_k-\eta g_k\bigr).
    \]
\EndFor
\State \textbf{Output:} $\{x_k\}_{k=0}^K$ (or the last iterate $x_K$).
\end{algorithmic}
\end{algorithm}

---------------------------------------------------------------------

\section*{Dry Run Example}
Consider the scalar problem $f(w)=(w-3)^2$ and $g\equiv0$.  
Thus $\Phi(w)=f(w)$, $\nabla f(w)=2(w-3)$.  
Take $\eta=0.1$, $w_0=0$, and run three iterations.

\begin{center}
\begin{tabular}{c|c|c|c}
$k$ & $w_k$ & $\nabla f(w_k)$ & $w_{k+1}=w_k-\eta\nabla f(w_k)$\\\hline
0 & $0$ & $2(0-3)=-6$ & $0-0.1(-6)=0.6$\\
1 & $0.6$ & $2(0.6-3)=-4.8$ & $0.6-0.1(-4.8)=1.08$\\
2 & $1.08$ & $2(1.08-3)=-3.84$ & $1.08-0.1(-3.84)=1.464$\\
\end{tabular}
\end{center}
Objective values:
\[
\Phi(w_0)=9,\;\Phi(w_1)=(0.6-3)^2=5.76,\;
\Phi(w_2)=(1.08-3)^2=3.6864,\;
\Phi(w_3)=(1.464-3)^2=2.365\,.
\]
The sequence monotonically decreases towards the minimiser $w^\star=3$.

---------------------------------------------------------------------

\section*{Assumptions}
\begin{assumption}[Smoothness of $f$]\label{ass:smooth}
$f:\mathbb{R}^d\to\mathbb{R}$ is differentiable and $L$‑Lipschitz gradient, i.e.
\[
\|\nabla f(x)-\nabla f(y)\|\le L\|x-y\|,\qquad\forall x,y\in\mathbb{R}^d .
\tag{5}
\]
Equivalently,
\[
f(y)\le f(x)+\langle\nabla f(x),y-x\rangle+\frac{L}{2}\|y-x\|^2 .
\tag{6}
\]
\end{assumption}

\begin{assumption}[Proximal regularity of $g$]\label{ass:prox}
$g:\mathbb{R}^d\to\mathbb{R}\cup\{+\infty\}$ is proper, lower‑semicontinuous and its proximal operator $\operatorname{prox}_{\eta g}$ is \emph{non‑expansive}:
\[
\bigl\|\operatorname{prox}_{\eta g}(u)-\operatorname{prox}_{\eta g}(v)\bigr\|
\le\|u-v\|,\qquad\forall u,v\in\mathbb{R}^d .
\tag{7}
\]
\end{assumption}

\begin{assumption}[Stepsize]\label{ass:stepsize}
The constant stepsize satisfies
\[
0<\eta\le\frac{1}{L}.
\tag{8}
\]
\end{assumption}

\begin{assumption}[Bounded below]\label{ass:bound}
The composite objective $\Phi$ is bounded from below:
\[
\Phi^\star\;\stackrel{\text{def}}{=}\;\inf_{x\in\mathbb{R}^d}\Phi(x)>-\infty .
\tag{9}
\]
\end{assumption}

---------------------------------------------------------------------

\section*{Main Convergence Theorem}
\begin{theorem}[Convergence of PGM‑NC]\label{thm:main}
Under Assumptions~\ref{ass:smooth}–\ref{ass:bound} and with a stepsize $\eta$ satisfying \eqref{ass:stepsize}, the iterates generated by \eqref{2} obey for every $K\ge1$
\[
\frac{1}{K}\sum_{k=0}^{K-1}\|G_{\eta}(x_k)\|^{2}
\;\le\;
\frac{2\bigl(\Phi(x_0)-\Phi^\star\bigr)}{\eta\,K}.
\tag{10}
\]
Consequently,
\[
\min_{0\le k<K}\|G_{\eta}(x_k)\|^{2}
\;\le\;
\frac{2\bigl(\Phi(x_0)-\Phi^\star\bigr)}{\eta\,K}
=O\!\left(\frac{1}{K}\right).
\tag{11}
\]
Thus any limit point of $\{x_k\}$ is a \emph{stationary point} of $\Phi$, i.e. $0\in\nabla f(\bar x)+\partial g(\bar x)$.
\end{theorem}

---------------------------------------------------------------------

\section*{Theoretical Properties}
\begin{itemize}
    \item \textbf{Stability:} The descent inequality \eqref{10} holds for any initial point $x_0$; the bound depends only on $\Phi(x_0)-\Phi^\star$.
    \item \textbf{Hyperparameter Sensitivity:} The stepsize must respect $\eta\le1/L$. Smaller $\eta$ yields a larger denominator in \eqref{10}, i.e. slower theoretical rate but increased robustness.
    \item \textbf{Comparison to Baselines:}
    \begin{enumerate}[label=(\alph*)]
        \item For pure gradient descent ($g\equiv0$) the bound reduces to the classical $O(1/K)$ rate for non‑convex smooth optimisation.
        \item Compared with stochastic variants, the deterministic PGM‑NC enjoys a deterministic $O(1/K)$ rate without variance terms.
    \end{enumerate}
    \item \textbf{Extension:} If $g$ is convex, the same proof gives a stronger result: $\Phi(x_{k})$ itself is non‑increasing.
\end{itemize}

---------------------------------------------------------------------

\section*{Discussion}
The theorem shows that the proximal gradient method retains the optimal $O(1/K)$ decrease of the squared gradient mapping even when $f$ is non‑convex, provided the smoothness constant $L$ is known to set $\eta\le1/L$. The proof crucially exploits the non‑expansiveness of the proximal operator \eqref{7} and the descent lemma \eqref{6}. No stochasticity is involved; hence the result is deterministic and holds pathwise.

---------------------------------------------------------------------

\section*{Preliminaries (Definitions and Lemmas)}
\begin{definition}[Gradient mapping]
For any $x\in\mathbb{R}^d$ and stepsize $\eta>0$ define
\[
G_{\eta}(x)\;\stackrel{\text{def}}{=}\;\frac{1}{\eta}\bigl(x-\operatorname{prox}_{\eta g}(x-\eta\nabla f(x))\bigr).
\tag{12}
\]
\end{definition}

\begin{lemma}[Non‑expansiveness of proximal operator]\label{lem:nonexp}
For any $u,v\in\mathbb{R}^d$,
\[
\bigl\|\operatorname{prox}_{\eta g}(u)-\operatorname{prox}_{\eta g}(v)\bigr\|
\le\|u-v\|.
\tag{13}
\]
\end{lemma}
\begin{proof}
Standard property of the proximal map of a proper l.s.c. function (see e.g. \cite{parikh2014proximal}). \qedhere
\end{proof}

\begin{lemma}[Descent lemma for $L$‑smooth $f$]\label{lem:descent}
For any $x,y\in\mathbb{R}^d$,
\[
f(y)\le f(x)+\langle\nabla f(x),y-x\rangle+\frac{L}{2}\|y-x\|^{2}.
\tag{14}
\]
\end{lemma}
\begin{proof}
Follows directly from Lipschitz continuity of $\nabla f$ (Assumption~\ref{ass:smooth}). \qedhere
\end{proof}

---------------------------------------------------------------------

\section*{Main Proof (Step‑by‑Step Derivation)}
We prove Theorem~\ref{thm:main}. All steps are labelled for easy reference.

\begin{proof}
Let $x_{k+1}$ be defined by \eqref{2}. Set $v_k\coloneqq x_k-\eta\nabla f(x_k)$. By definition of the proximal operator,
\[
x_{k+1}\;=\;\arg\min_{u}\Bigl\{g(u)+\frac{1}{2\eta}\|u-v_k\|^{2}\Bigr\}.
\tag{15}
\]
The optimality condition (first‑order condition for the subdifferential) yields
\[
0\;\in\;\partial g(x_{k+1})+\frac{1}{\eta}\bigl(x_{k+1}-v_k\bigr)
\;\Longleftrightarrow\;
\frac{1}{\eta}\bigl(v_k-x_{k+1}\bigr)\;\in\;\partial g(x_{k+1}).
\tag{16}
\]
Recall $v_k=x_k-\eta\nabla f(x_k)$. Substituting into \eqref{16} gives
\[
\frac{1}{\eta}\bigl(x_k-x_{k+1}\bigr)-\nabla f(x_k)\;\in\;\partial g(x_{k+1}).
\tag{17}
\]
Define the gradient mapping $G_{\eta}(x_k)$ as in \eqref{12}; then \eqref{17} can be rewritten as
\[
-\bigl(G_{\eta}(x_k)+\nabla f(x_k)\bigr)\;\in\;\partial g(x_{k+1}).
\tag{18}
\]

\medskip
\noindent\textbf{[Step 1] Descent of the smooth part.}  
Apply Lemma~\ref{lem:descent} with $x=x_k$, $y=x_{k+1}$:
\[
f(x_{k+1})\le f(x_k)+\langle\nabla f(x_k),x_{k+1}-x_k\rangle+\frac{L}{2}\|x_{k+1}-x_k\|^{2}.
\tag{19}
\]

\noindent\textbf{[Step 2] Descent of the nonsmooth part via proximal optimality.}  
From the definition \eqref{15} we have for any $u$,
\[
g(x_{k+1})+\frac{1}{2\eta}\|x_{k+1}-v_k\|^{2}
\;\le\;
g(u)+\frac{1}{2\eta}\|u-v_k\|^{2}.
\tag{20}
\]
Choose $u=x_k$. Using $v_k=x_k-\eta\nabla f(x_k)$, the right‑hand side becomes
\[
g(x_k)+\frac{1}{2\eta}\|x_k-(x_k-\eta\nabla f(x_k))\|^{2}
= g(x_k)+\frac{\eta}{2}\|\nabla f(x_k)\|^{2}.
\tag{21}
\]
The left‑hand side expands to
\[
g(x_{k+1})+\frac{1}{2\eta}\|x_{k+1}-x_k+\eta\nabla f(x_k)\|^{2}
= g(x_{k+1})+\frac{1}{2\eta}\bigl(\|x_{k+1}-x_k\|^{2}
+2\eta\langle x_{k+1}-x_k,\nabla f(x_k)\rangle
+\eta^{2}\|\nabla f(x_k)\|^{2}\bigr).
\tag{22}
\]
Inequality \eqref{20} together with \eqref{21}–\eqref{22} yields
\[
g(x_{k+1})\le g(x_k)
-\langle\nabla f(x_k),x_{k+1}-x_k\rangle
-\frac{1}{2\eta}\|x_{k+1}-x_k\|^{2}.
\tag{23}
\]

\noindent\textbf{[Step 3] Combine smooth and nonsmooth descents.}  
Add \eqref{19} and \eqref{23}:
\[
\begin{aligned}
\Phi(x_{k+1})
&=f(x_{k+1})+g(x_{k+1})\\
&\le f(x_k)+g(x_k)
-\frac{1}{2\eta}\|x_{k+1}-x_k\|^{2}
+\frac{L}{2}\|x_{k+1}-x_k\|^{2}.
\end{aligned}
\tag{24}
\]
Since $\eta\le 1/L$ (Assumption~\ref{ass:stepsize}), we have $L\eta\le1$ and therefore
\[
\frac{L}{2}\|x_{k+1}-x_k\|^{2}
\le\frac{1}{2\eta}\|x_{k+1}-x_k\|^{2}.
\tag{25}
\]
Plugging \eqref{25} into \eqref{24} gives the key descent inequality
\[
\Phi(x_{k+1})\le\Phi(x_k)-\frac{1}{2\eta}\|x_{k+1}-x_k\|^{2}.
\tag{26}
\]

\noindent\textbf{[Step 4] Relate the step norm to the gradient mapping.}  
Recall $G_{\eta}(x_k)=\frac{1}{\eta}(x_k-x_{k+1})$, thus
\[
\|x_{k+1}-x_k\|^{2}=\eta^{2}\|G_{\eta}(x_k)\|^{2}.
\tag{27}
\]
Insert \eqref{27} into \eqref{26}:
\[
\Phi(x_{k+1})\le\Phi(x_k)-\frac{\eta}{2}\|G_{\eta}(x_k)\|^{2}.
\tag{28}
\]

\noindent\textbf{[Step 5] Summation over iterations.}  
Sum \eqref{28} for $k=0,\dots,K-1$:
\[
\Phi(x_K)\le\Phi(x_0)-\frac{\eta}{2}\sum_{k=0}^{K-1}\|G_{\eta}(x_k)\|^{2}.
\tag{29}
\]
Rearrange and use $\Phi(x_K)\ge\Phi^\star$ (Assumption~\ref{ass:bound}):
\[
\frac{1}{K}\sum_{k=0}^{K-1}\|G_{\eta}(x_k)\|^{2}
\le\frac{2\bigl(\Phi(x_0)-\Phi^\star\bigr)}{\eta\,K}.
\tag{30}
\]
Equation \eqref{30} is exactly the bound \eqref{10}. The inequality for the minimum follows immediately:
\[
\min_{0\le k<K}\|G_{\eta}(x_k)\|^{2}
\le\frac{1}{K}\sum_{k=0}^{K-1}\|G_{\eta}(x_k)\|^{2}
\le\frac{2\bigl(\Phi(x_0)-\Phi^\star\bigr)}{\eta\,K}.
\tag{31}
\]
Thus the theorem holds. \qedhere
\end{proof}

---------------------------------------------------------------------

\section*{Rate Derivation (With Constants)}
From \eqref{30} we obtain the explicit constant
\[
C\;=\;\frac{2\bigl(\Phi(x_0)-\Phi^\star\bigr)}{\eta},
\qquad
\frac{1}{K}\sum_{k=0}^{K-1}\|G_{\eta}(x_k)\|^{2}\le\frac{C}{K}.
\tag{32}
\]
If the stepsize is chosen as the maximal admissible value $\eta=1/L$, then
\[
C = 2L\bigl(\Phi(x_0)-\Phi^\star\bigr).
\tag{33}
\]

---------------------------------------------------------------------

\section*{Special Cases}
\begin{enumerate}[label=(\alph*)]
    \item \textbf{Pure gradient descent ($g\equiv0$).}  
    The proximal step reduces to $x_{k+1}=x_k-\eta\nabla f(x_k)$.  
    The gradient mapping becomes $G_{\eta}(x_k)=\nabla f(x_k)$, and \eqref{30} gives
    \[
    \frac{1}{K}\sum_{k=0}^{K-1}\|\nabla f(x_k)\|^{2}\le\frac{2\bigl(f(x_0)-f^\star\bigr)}{\eta K},
    \]
    the classical $O(1/K)$ rate for non‑convex smooth optimisation.

    \item \textbf{Convex $g$.}  
    When $g$ is convex, the optimality condition \eqref{16} yields the stronger inequality
    \[
    \Phi(x_{k+1})\le\Phi(x_k)-\frac{1}{2\eta}\|x_{k+1}-x_k\|^{2},
    \]
    and consequently $\{\Phi(x_k)\}$ is monotone decreasing.

    \item \textbf{Exact line‑search stepsize.}  
    If at each iteration we select $\eta_k$ via the Armijo rule ensuring
    \[
    \Phi(x_{k+1})\le\Phi(x_k)-\frac{\eta_k}{2}\|G_{\eta_k}(x_k)\|^{2},
    \]
    the same summation argument provides a bound with $\eta_{\min}=\min_k\eta_k$ in place of $\eta$.
\end{enumerate}

---------------------------------------------------------------------

\begin{thebibliography}{9}
\bibitem{parikh2014proximal}
N.~Parikh and S.~Boyd, ``Proximal algorithms,'' \emph{Foundations and Trends in Optimization}, vol.~1, no.~3, pp. 127–239, 2014.
\end{thebibliography}

\end{document}